% GENERATED FILE. Edit canonical lessons, then run npm run generate:books.
% canonical-source-hash: fnv1a64:532bdbbb36584870
% canonical-lessons: FR-C387-proces, FR-C387-defense, FR-C387-verdict, FR-C387-coupable, FR-C387-innocent

\chapter{Au tribunal}
\label{ch:fr-au-tribunal}
\hlchaptermodality{387}

\begin{chapteropening}
\textbf{By the end of this chapter:} \emph{I can describe a trial: the defence, the verdict, and whether someone is guilty or innocent.}

Complete the last lesson of chapter 387: I can describe a trial: the defence, the verdict, and whether someone is guilty or innocent.
\end{chapteropening}

\section[le procès]{le procès — the trial}

\label{lesson:FR-C387-proces}



\begin{quote}
Before the new one: say the French for the complaint, then the French for the misdemeanour.
\end{quote}

\subsection*{You'll want to know: le procès}
\textbf{le procès} — \textquotedblleft{}the trial\textquotedblright{}.

\begin{quote}
Le procès commence la semaine prochaine.
\end{quote}

\begin{quote}
\emph{The trial starts next week.}
\end{quote}

\emph{Faire un procès à quelqu'un} — to sue someone.

\begin{cousinweb}
\emph{procès} comes from Latin \emph{prōcessus} \textquotedblleft{}a going forward, progress\textquotedblright{} (EtymDB; Lewis \& Short), from \emph{prōcēdere} \textquotedblleft{}to go forward\textquotedblright{}, built on \emph{cēdere} \textquotedblleft{}to go\textquotedblright{}.
\end{cousinweb}

\begin{grammarlens}[title={What you've built}]
One more word for telling, explaining and describing at B1.
\end{grammarlens}

\subsection*{Your turn}
\begin{itemize}
\raggedright
  \item \emph{Say it:} \emph{le procès}
  \item \emph{Say it:} \emph{le procès}, once more
  \item \emph{Say it:} a sentence of your own with \emph{le procès}
  \item \emph{Recall:} say the French for the complaint, then the French for the misdemeanour, then say \emph{le procès} again
\end{itemize}

\begin{tcolorbox}[breakable,colback=teal!4,colframe=teal!35!black,title={Before you move on}]
What does le procès mean? (\textquotedblleft{}The trial\textquotedblright{}.) Say it once more.
\end{tcolorbox}
\section[la défense]{la défense — the defence}

\label{lesson:FR-C387-defense}



\begin{quote}
Before the new one: say the French for the misdemeanour, then the French for the trial.
\end{quote}

\subsection*{You'll want to know: la défense}
\textbf{la défense} — \textquotedblleft{}the defence\textquotedblright{}.

\begin{quote}
L'avocat de la défense a parlé pendant une heure.
\end{quote}

\begin{quote}
\emph{The defence lawyer spoke for an hour.}
\end{quote}

\emph{Défense de fumer} — no smoking.

\begin{cousinweb}
\emph{défense} comes from Late Latin \emph{dēfēnsa} \textquotedblleft{}protection\textquotedblright{} (EtymDB), from \emph{dēfendere} \textquotedblleft{}to ward off, to defend\textquotedblright{} (Lewis \& Short).
\end{cousinweb}

\begin{grammarlens}[title={What you've built}]
One more word for telling, explaining and describing at B1.
\end{grammarlens}

\subsection*{Your turn}
\begin{itemize}
\raggedright
  \item \emph{Say it:} \emph{la défense}
  \item \emph{Say it:} \emph{la défense}, once more
  \item \emph{Say it:} a sentence of your own with \emph{la défense}
  \item \emph{Recall:} say the French for the misdemeanour, then the French for the trial, then say \emph{la défense} again
\end{itemize}

\begin{tcolorbox}[breakable,colback=teal!4,colframe=teal!35!black,title={Before you move on}]
What does la défense mean? (\textquotedblleft{}The defence\textquotedblright{}.) Say it once more.
\end{tcolorbox}
\section[le verdict]{le verdict — the verdict}

\label{lesson:FR-C387-verdict}



\begin{quote}
Before the new one: say the French for the trial, then the French for the defence.
\end{quote}

\subsection*{You'll want to know: le verdict}
\textbf{le verdict} — \textquotedblleft{}the verdict\textquotedblright{}.

\begin{quote}
Le verdict tombe demain.
\end{quote}

\begin{quote}
\emph{The verdict comes tomorrow.}
\end{quote}

The final \emph{ct} is said.

\begin{cousinweb}
\emph{verdict} was borrowed from English \emph{verdict}, from Medieval Latin \emph{vērēdictum} \textquotedblleft{}a thing truly said\textquotedblright{}, built from \emph{vērē} \textquotedblleft{}truly\textquotedblright{} and \emph{dictum} \textquotedblleft{}said\textquotedblright{} (Lewis \& Short).
\end{cousinweb}

\begin{grammarlens}[title={What you've built}]
One more word for telling, explaining and describing at B1.
\end{grammarlens}

\subsection*{Your turn}
\begin{itemize}
\raggedright
  \item \emph{Say it:} \emph{le verdict}
  \item \emph{Say it:} \emph{le verdict}, once more
  \item \emph{Say it:} a sentence of your own with \emph{le verdict}
  \item \emph{Recall:} say the French for the trial, then the French for the defence, then say \emph{le verdict} again
\end{itemize}

\begin{tcolorbox}[breakable,colback=teal!4,colframe=teal!35!black,title={Before you move on}]
What does le verdict mean? (\textquotedblleft{}The verdict\textquotedblright{}.) Say it once more.
\end{tcolorbox}
\section[coupable]{coupable — guilty}

\label{lesson:FR-C387-coupable}



\begin{quote}
Before the new one: say the French for the defence, then the French for the verdict.
\end{quote}

\subsection*{You'll want to know: coupable}
\textbf{coupable} — \textquotedblleft{}guilty\textquotedblright{}.

\begin{quote}
Pour le juge, il est coupable.
\end{quote}

\begin{quote}
\emph{For the judge, he is guilty.}
\end{quote}

Also a noun: \emph{le coupable} — the culprit.

\begin{cousinweb}
\emph{coupable} is inherited from Latin \emph{culpābilis} \textquotedblleft{}blameworthy\textquotedblright{} (EtymDB), from \emph{culpa} \textquotedblleft{}fault, blame\textquotedblright{} (Lewis \& Short).
\end{cousinweb}

\begin{grammarlens}[title={What you've built}]
One more word for telling, explaining and describing at B1.
\end{grammarlens}

\subsection*{Your turn}
\begin{itemize}
\raggedright
  \item \emph{Say it:} \emph{coupable}
  \item \emph{Say it:} \emph{coupable}, once more
  \item \emph{Say it:} a sentence of your own with \emph{coupable}
  \item \emph{Recall:} say the French for the defence, then the French for the verdict, then say \emph{coupable} again
\end{itemize}

\begin{tcolorbox}[breakable,colback=teal!4,colframe=teal!35!black,title={Before you move on}]
What does coupable mean? (\textquotedblleft{}Guilty\textquotedblright{}.) Say it once more.
\end{tcolorbox}
\section[innocent]{innocent — innocent}

\label{lesson:FR-C387-innocent}



\begin{quote}
Before the new one: say the French for the verdict, then the French for guilty.
\end{quote}

\subsection*{You'll want to know: innocent}
\textbf{innocent} — \textquotedblleft{}innocent\textquotedblright{}.

\begin{quote}
Il dit qu'il est innocent.
\end{quote}

\begin{quote}
\emph{He says he is innocent.}
\end{quote}

Feminine \emph{innocente}.

\begin{cousinweb}
\emph{innocent} is a learned borrowing from Latin \emph{innocēns} \textquotedblleft{}harmless, blameless\textquotedblright{} (EtymDB; Lewis \& Short), built from \emph{in-} \textquotedblleft{}not\textquotedblright{} and \emph{nocēre} \textquotedblleft{}to harm\textquotedblright{}.
\end{cousinweb}

\begin{grammarlens}[title={What you've built}]
One more word for telling, explaining and describing at B1.
\end{grammarlens}

\subsection*{Your turn}
\begin{itemize}
\raggedright
  \item \emph{Say it:} \emph{innocent}
  \item \emph{Say it:} \emph{innocent}, once more
  \item \emph{Say it:} a sentence of your own with \emph{innocent}
  \item \emph{Recall:} say the French for the verdict, then the French for guilty, then say \emph{innocent} again
\end{itemize}

\begin{tcolorbox}[breakable,colback=teal!4,colframe=teal!35!black,title={Before you move on}]
What does innocent mean? (\textquotedblleft{}Innocent\textquotedblright{}.) Say it once more.
\end{tcolorbox}
